% R008 — Guide to Cultivating Complex Analysis (Bahasa Indonesia)
% Reader-facing translation unit: source/ca-v1.9/ca.tex lines 1265–1288.
% Source release: v1.9 (commit a4d25d9ba37a486a5a020219eb8bd4e6602fd14c).
% Translation provenance: OpenAI Codex gpt-5.6-sol, Ultra, acting on the user's request.
% Preserve source equations, labels, and figure paths.

\subsection{Koordinat polar}

Karena bilangan kompleks dapat dipandang sebagai titik-titik pada bidang, kita dapat menggunakan
\emph{\myindex{koordinat polar}}
untuk merepresentasikan bilangan kompleks.
Yaitu, $x = r \cos \theta$ dan $y= r \sin \theta$.
Kita menuliskan $z$ dalam \emph{\myindex{bentuk polar}} sebagai
\begin{equation*}
z = x+iy = r e^{i \theta} .
\end{equation*}
Di sini, $r = \sabs{z} = \sqrt{x^2+y^2}$ adalah modulus, dan
$\theta$ adalah sudut yang dibentuk oleh $x+iy$ dengan sumbu real positif (sumbu-$x$).
$\theta$ disebut \emph{\myindex{argumen}}. Lihat
\figureref{fig:polarcoords}.
Alasan untuk notasi ini adalah
rumus Euler, sehingga
\begin{equation*}
z = r e^{i\theta} = r\cos \theta + i r\sin \theta  = x+iy .
\end{equation*}

\begin{myfig}
\subimport*{figures/}{polarcoords.pdf_t}
\caption{Koordinat polar.\label{fig:polarcoords}}
\end{myfig}
